% Part: first-order-logic
% Chapter: completeness
% Section: outline

\documentclass[../../../include/open-logic-section]{subfiles}

\begin{document}

\iftag{FOL}
      {\olfileid[hi]{fol}{com}{out}}
      {\olfileid[hi]{pl}{com}{out}}

\olsection{प्रमाण की रूपरेखा}

पूर्णता प्रमेय का प्रमाण कुछ जटिल है और उसे पहली बार पढ़ते समय रास्ता भटकना
आसान है। इसलिए पहले प्रमाण की रूपरेखा देख लें। पहला चरण दृष्टिकोण का ऐसा
परिवर्तन है जो हमें प्रमाण तक पहुँचने का मार्ग दिखाता है। यदि पूर्णता को
``जब कभी $\Gamma
\Entails !A$ हो, तब $\Gamma \Proves !A$ होता है'' के रूप में सोचें, तो कोई
युक्ति सूझना भी कठिन हो सकता है: $\Gamma \Proves !A$ दिखाने के लिए हमें
!!a{derivation} खोजनी है, और यह नहीं दिखता कि $\Gamma \Entails !A$ वाली
परिकल्पना इसमें किसी प्रकार हमारी सहायता करती है। कुछ प्रमाण-प्रणालियों के
लिए !!a{derivation} का सीधे निर्माण संभव है, पर हम थोड़ा भिन्न मार्ग लेंगे।
आवश्यक दृष्टिकोण-परिवर्तन यह है: पूर्णता को इस रूप में भी कहा जा सकता है:
``यदि $\Gamma$ संगत है, तो वह संतुष्टनीय है।'' संभव है कि हम~$\Gamma$ में
निहित सूचना को उसकी संगतता की परिकल्पना के साथ लेकर ऐसी
\iftag{FOL}{!!a{structure}}{!!a{valuation}} बना सकें जो~$\Gamma$ के प्रत्येक
\iftag{FOL}{!!{sentence}}{!!{formula}} को संतुष्ट करे। आखिर हमें मालूम है कि
हम किस प्रकार की \iftag{FOL}{!!{structure}}{!!{valuation}} खोज रहे हैं: ऐसी
जो ठीक वैसी हो जैसी $\Gamma$ उसका वर्णन करता है!{}

यदि $\Gamma$ में केवल \iftag{FOL}{परमाण्विक
  !!{sentence}s}{!!{propositional variable}s} हों, तो उसका प्रतिरूप बनाना
सरल है।\iftag{FOL}{ मान लें कि सभी परमाण्विक !!{sentence}s का रूप
  $\Atom{P}{a_1,\dots,a_n}$ है, जहाँ $a_i$ सभी !!{constant}s हैं।}{} हमें बस
\iftag{FOL}{%
  !!a{domain}~$\Domain{M}$ और~$P$ के लिए ऐसा अर्थ-निर्धारण चुनना है कि
  $\Sat{M}{\Atom{P}{a_1,\ldots,a_n}}$ हो। यह अधिक कठिन नहीं है:
  $\Domain{M} = \Nat$ रखें, $\Assign{\Obj c_i}{M} = i$ रखें, और प्रत्येक
  $\Atom{P}{a_1,\ldots,a_n} \in \Gamma$ के लिए टपल $\tuple{k_1,
    \dots, k_n}$ को $\Assign{P}{M}$ में रखें, जहाँ $k_i$ अचर
  चिह्न~$a_i$ का सूचकांक है (अर्थात् $a_i \ident \Obj c_{k_i}$)।}{%
  ऐसी !!a{valuation}~$\pAssign{v}$ चुननी है कि $\pSat{v}{p}$ सभी $p \in
  \Gamma$ के लिए हो। इसके लिए $\pAssign{v}(p) = \True$ तभी
  रखें जब $p \in \Gamma$ हो।}

अब मान लें कि $\Gamma$ में कोई !!{formula}~$\lnot !B$ है, जहाँ $!B$
परमाण्विक है। हमें आशंका हो सकती है कि
$\iftag{FOL}{\Struct{M}}{\pAssign{v}}$ का निर्माण $\lnot !B$ को सत्य बनाने की
संभावना में बाधा डालेगा। किंतु यहीं~$\Gamma$ की संगतता काम आती है: यदि
$\lnot !B \in \Gamma$ है, तो $!B \notin
\Gamma$ होगा, अन्यथा $\Gamma$ असंगत होता। और यदि $!B \notin
\Gamma$ है, तो हमारे
$\iftag{FOL}{\Struct{M}}{\pAssign{v}}$ के निर्माण के अनुसार
$\iftag{FOL}{\Sat/{M}}{\pSat/{v}}{!B}$ है, इसलिए
$\iftag{FOL}{\Sat{M}}{\pSat{v}}{\lnot !B}$। यहाँ तक सब ठीक है।

यदि $\Gamma$ में जटिल, गैर-परमाण्विक सूत्र हों तो क्या होगा? मान लें कि
उसमें $!A \land !B$ है। इसे सत्य बनाने के लिए हमें ऐसे आगे बढ़ना चाहिए मानो
$!A$ और $!B$ दोनों~$\Gamma$ में हों। और यदि $!A \lor !B \in \Gamma$ है,
तो हमें उनमें से कम-से-कम एक को सत्य बनाना होगा, अर्थात् ऐसे आगे बढ़ना होगा
मानो उनमें से कोई एक~$\Gamma$ में हो।

इससे निम्न विचार सूझता है: हम~$\Gamma$ में अतिरिक्त !!{formula}s इस प्रकार
जोड़ें कि (a)~परिणामी समुच्चय संगत बना रहे, (b)~प्रत्येक संभावित परमाण्विक
!!{sentence}~$!A$ के लिए या तो $!A$ परिणामी समुच्चय में हो या $\lnot !A$ हो,
और (c)~जब भी $!A \land !B$ समुच्चय में हो, तब $!A$ और $!B$ दोनों भी उसमें
हों; यदि $!A \lor !B$ समुच्चय में हो, तो $!A$ या $!B$ में से कम-से-कम एक
भी उसमें हो; इत्यादि। हम यह प्रक्रिया (संभवतः अनंत काल तक) जारी रखते हैं।
इस प्रकार जोड़े गए सभी !!{formula}s के समुच्चय को~$\Gamma^*$ कहें। तब ऊपर
का हमारा निर्माण हमें ऐसी \iftag{FOL}{!!a{structure}~$\Struct{M}$}{!!a{valuation}~$\pAssign{v}$}
देगा जिसके बारे में हम आगमन से सिद्ध कर सकेंगे कि वह~$\Gamma^*$ के सभी
वाक्यों को संतुष्ट करती है, और इस कारण~$\Gamma$ के सभी वाक्यों को भी,
क्योंकि $\Gamma \subseteq \Gamma^*$। पता चलता है कि (a) और~(b) की गारंटी
ही पर्याप्त है। वाक्यों का ऐसा समुच्चय जिसके लिए (b) सत्य हो,
\emph{पूर्ण} कहलाता है। इसलिए हमारा कार्य संगत समुच्चय~$\Gamma$ को संगत और
पूर्ण समुच्चय~$\Gamma^*$ तक विस्तारित करना होगा।

\iftag{FOL}{%
इस योजना में एक अड़चन है: यदि $\lexists[x][!A(x)] \in \Gamma$ है, तो हम
चाहेंगे कि कोई !!{constant}~$c$ चुनकर इस प्रक्रिया में $!A(c)$ जोड़ सकें।
पर हम कैसे जानें कि ऐसा करना सदा संभव है? संभव है कि हमारी भाषा में केवल
कुछ !!{constant}s हों और उनमें से प्रत्येक के लिए $\lnot !A(c) \in \Gamma$
हो। हम $!A(c)$ भी नहीं जोड़ सकते, क्योंकि उससे समुच्चय असंगत हो जाएगा और
हमें मालूम नहीं होगा कि $\Struct{M}$ को $!A(c)$ सत्य बनाना है या
$\lnot !A(c)$। यह भी संभव है कि $\Gamma$ में केवल ऐसी भाषा के वाक्य हों
जिसमें एक भी अचर चिह्न न हो (उदाहरणतः समुच्चय-सिद्धांत की भाषा)।

इस समस्या का समाधान यह है कि आरंभ में ही अनंततः अनेक अचर और ऐसे वाक्य जोड़
दिए जाएँ जो उन्हें उचित ढंग से परिमाणकों से जोड़ते हों। (निस्संदेह, हमें
सत्यापित करना होगा कि इससे असंगति उत्पन्न नहीं हो सकती।)

हमारा मूल निर्माण उस समय सुचारु है जब परमाण्विक वाक्यों में केवल
!!{constant}s हों। पर भाषा में !!{function}s भी हो सकते हैं। उस दशा में इन
!!{function}s को~$\Nat$ पर उचित फलन देना, ताकि सब कुछ काम करे, कठिन हो सकता
है। इसलिए एक और युक्ति अपनाएँ: $i$ से $\Obj c_i$ की व्याख्या करने के बदले,
स्वयं !!{constant}s के समुच्चय को प्रांत मान लें। तब
$\Struct M$ प्रत्येक !!{constant} को उसी को निर्दिष्ट कर सकती है:
$\Assign{\Obj c_i}{M} = \Obj c_i$। किंतु यहीं क्यों रुकें: $\Domain{M}$ को
भाषा के सभी \emph{पदों} का समुच्चय मान लें!{} ऐसा करने पर !!{function}s को
फलन (जो पदों को निविष्टि के रूप में लेकर पदों को मान के रूप में देते हैं)
निर्दिष्ट करने का एक प्रत्यक्ष उपाय है: !!{function}~$\Obj f^n_i$ को वह फलन
दें जो $n$ पदों $t_1$, \dots,~$t_n$ को निविष्टि के रूप में लेकर
$\Obj f^n_i(t_1, \dots, t_n)$ पद को मान के रूप में देता है।

पहेली का अंतिम टुकड़ा यह है कि~$\eq$ का क्या किया जाए। !!{predicate}~$\eq$
की व्याख्या नियत है: $\Sat{M}{\eq[t][t']}$ तभी जब $\Value{t}{M} = \Value{t'}{M}$
हो। अब यदि हम व्यवस्था ऐसी करें कि पद~$t$ का !!{value} स्वयं
$t$ हो, तो यह !!{structure} $\eq[t][t']$ रूप के \emph{किसी भी} वाक्य को तब
तक सत्य नहीं बनाएगी जब तक $t$ और $t'$ बिल्कुल वही पद न हों। यह निस्संदेह
एक समस्या है, क्योंकि !!{function}s वाली भाषा के लगभग प्रत्येक रोचक
सिद्धांत में $\eq[t][t']$ रूप के ऐसे वाक्य प्रमेय होंगे जिनमें $t$ और $t'$
एक ही पद नहीं हैं (उदाहरणतः अंकगणित के सिद्धांतों में:
$\eq[(\Obj 0+ \Obj 0)][\Obj 0]$)। इस समस्या के समाधान के लिए हम~$\Struct M$
का प्रांत बदलते हैं:~$\Domain{M}$ की वस्तुओं के रूप में पदों का उपयोग करने
के बजाय हम पदों के समुच्चयों का उपयोग करते हैं, और प्रत्येक समुच्चय में वे
सभी पद होते हैं जिन्हें~$\Gamma$ के वाक्य बराबर होने को बाध्य करते हैं।
अतः, उदाहरण के लिए, यदि $\Gamma$ अंकगणित का सिद्धांत है, तो इनमें से एक
समुच्चय में ये सब होंगे: $\Obj 0$, $(\Obj 0 + \Obj 0)$, $(\Obj
0 \times \Obj 0)$, इत्यादि। इसी समुच्चय को हम $\Obj 0$ के लिए निर्दिष्ट
करेंगे, और पता चलेगा कि यही समुच्चय इसके सभी पदों का भी मान है, जैसे
$(\Obj 0 + \Obj 0)$ का। इसलिए वाक्य
$\eq[(\Obj 0+ \Obj 0)][\Obj 0]$ इस संशोधित !!{structure} में सत्य होगा।}{}

अब देखें कि हम ठीक क्या करेंगे। पहले हम !!{complete} संगत समुच्चयों के गुणों
की जाँच करेंगे। विशेषतः हम सिद्ध करेंगे कि !!a{complete} संगत समुच्चय में
$!A \land !B$ तभी होता है जब उसमें $!A$ और~$!B$ दोनों हों, और $!A \lor !B$
तभी होता है जब उनमें से कम-से-कम एक उसमें हो; इत्यादि
(\olref[ccs]{prop:ccs})।\iftag{FOL}{ इसके बाद हम वाक्यों के ``संतृप्त''
  समुच्चयों को परिभाषित करेंगे और उनके गुणों की जाँच करेंगे। संतृप्त समुच्चय
  वह है जिसमें ऐसे सशर्त वाक्य होते हैं जो प्रत्येक परिमाणित !!{sentence}
  को उसके निदर्शों से जोड़ते हैं (\olref[hen]{defn:henkin-exp})। हम दिखाएँगे
  कि किसी भी संगत समुच्चय~$\Gamma$ को सदा किसी संतृप्त समुच्चय~$\Gamma'$ तक
  विस्तारित किया जा सकता है (\olref[hen]{lem:henkin})। यदि कोई समुच्चय संगत,
  संतृप्त और !!{complete} है, तो उसमें यह गुण भी होता है कि वह
  \iftag{prvEx}{$\lexists[x][!A(x)]$ को तभी समाहित करता है जब वह $!A(t)$ को
    किसी बंद पद~$t$ के लिए समाहित करे}{}\iftag{defEx,defAll}{}{ और
  }\iftag{prvAll}{$\lforall[x][!A(x)]$ को तभी समाहित करता है जब वह $!A(t)$ को
    प्रत्येक बंद पद~$t$ के लिए समाहित करे}{}
  (\olref[hen]{prop:saturated-instances})।}{} फिर हम
\iftag{FOL}{संतृप्त}{} संगत समुच्चय~$\iftag{FOL}{\Gamma'}{\Gamma}$ लेकर
दिखाएँगे कि उसे \iftag{FOL}{संतृप्त, संगत और}{संगत और} !!{complete}
समुच्चय~$\Gamma^*$ तक विस्तारित किया जा सकता है
(\olref[lin]{lem:lindenbaum})। इसी समुच्चय $\Gamma^*$ से हम अपना
\iftag{FOL}{पद-प्रतिरूप~$\Struct M(\Gamma^*)$}{!!{valuation}~$\pAssign v(\Gamma^*)$}
परिभाषित करेंगे। \iftag{FOL}{पद-प्रतिरूप का प्रांत बंद पदों का समुच्चय है
  और उसके !!{predicate}s की व्याख्या उन परमाण्विक
  !!{sentence}s}{सत्य-मूल्यांकन का निर्धारण उन !!{propositional
    variable}s} से होता है जो~$\Gamma^*$ में हैं
(\olref[mod]{defn:termmodel})। हम
\iftag{FOL}{संतृप्त,}{} पूर्ण संगत समुच्चयों के गुणों से दिखाएँगे कि वास्तव
में $\iftag{FOL}{\Sat{M(\Gamma^*)}}{\pSat{v(\Gamma^*)}}{!A}$ तभी जब $!A \in
\Gamma^*$ हो (\olref[mod]{lem:truth}), और अतः विशेषतः
$\iftag{FOL}{\Sat{M(\Gamma^*)}}{\pSat{v(\Gamma^*)}}{\Gamma}$।
\iftag{FOL}{अंत में, हम विचार करेंगे कि यदि $\Gamma$ में~$\eq$ भी हो तो
  पद-प्रतिरूप कैसे परिभाषित किया जाए
  (\olref[ide]{defn:term-model-factor}), और दिखाएँगे कि वह~$\Gamma^*$ को
  संतुष्ट करता है (\olref[ide]{lem:truth})।}{}

\end{document}
